% The design language (NNB text): included by nn_builder_manual.tex. Every example is a file in examples/ that the tests build
% (crates/core/tests/text.rs, manual_examples_build), so what is shown here works as printed.
\section{Writing designs as text: the NNB language}
\label{sec:text}

Every design can also be written as text, in the \ui{Text} tab. The text and the canvas are two views of the same design and are kept in
step both ways: typing changes the canvas a moment after the text is valid, and a change on the canvas changes the text. This chapter
teaches the language from the first line to functions and loops, then lists every block and every rule.

Why write text at all? A network of fifty blocks is fifty short lines, typed in a minute; every number of every block is visible at once;
repeated parts are written once (\texttt{repeat 40 \{ \dots\ \}}); and the text can be copied into a paper, an e-mail or a chat and pasted
back unchanged (\ui{Copy}, or \ui{File} $\rightarrow$ \ui{Export Design\dots} as a \file{.nnb} file).

\subsection{A first design, line by line}
The smallest useful design is a chain of blocks. LeNet-5 (Section~\ref{sec:demos}) for the $1\times28\times28$ MNIST digits:

\VerbatimInput[fontsize=\small,frame=single,numbers=left,numbersep=6pt,label=\file{lenet5.nnb}]{examples/lenet5.nnb}

\begin{itemize}[nosep]
  \item \texttt{net "LeNet-5"} names the design (optional).
  \item Each further line is one block: a \textbf{label} (\texttt{1}, \texttt{2}, \dots), the \textbf{block} (\texttt{conv},
        \texttt{avgpool}, \dots) and its \textbf{settings} in any order. Left-out settings take the same defaults as the canvas (a
        \texttt{conv} without a kernel is $3\times3$, without an activation it uses ReLU, \dots).
  \item A block takes its input from \textbf{the line above} unless it says otherwise.
  \item \texttt{\#} starts a comment, to the end of the line. Blank lines and indentation do not matter.
  \item The first block is \texttt{input}: its shape comes from the Dataset tab, not from the text. \texttt{output} gets as many
        classes as the dataset has.
\end{itemize}
Typed into the Text tab, this draws the nine blocks on the canvas, each with its label in its header, and the shapes and parameters are
those of Section~\ref{sec:math}: 61\,706 parameters.

\subsection{Settings: words and numbers}
Settings are written the way you would say them. For a convolution, the first plain number is the number of filters, \texttt{5x5} is
the kernel, and words set the rest:

\begin{verbatim}
2 conv 64 3x3                      # 64 filters, 3×3, stride 1, padding "same", ReLU
3 conv 128 3x3 stride 2 pad 1 gelu # stride 2, padding 1, GELU
4 conv 128 1x1 valid none nobias   # no padding, no activation, no bias
5 conv 64 3x3 dil 2                # dilation 2
6 conv filters=64 kernel=7x7 stride=2 padding=3   # any setting as name=value
\end{verbatim}
Every setting also has a \emph{name}\texttt{=}\emph{value} form (the names are those of the inspector), useful when a value is computed
(\texttt{kernel=s+2}, below). A word that a block does not know is an error, never ignored: \texttt{conv 64 stide 2} says
``\texttt{'stide' means nothing for conv (did you mean 'stride'?)}''. A value outside its range (a kernel of 40, 100\,000 filters) is an
error on its block.

\subsection{Connections: \texttt{<-}}
A block with a different input, or with several (Add, Concat, Multiply), names them after \texttt{<-}:

\begin{verbatim}
1 input
2 conv 64 3x3
3 conv 64 3x3 none
4 add <- 2, 3          # residual: the output of block 2 plus that of block 3
5 relu
6 conv 32 1x1 <- 2     # a side branch from block 2 (not from the line above)
7 concat <- 5, 6       # channels joined in this order: 5's first, then 6's
\end{verbatim}
A block may only use blocks above it (a network computes each block once, from the blocks before it), so the text can never contain a
loop. \texttt{<- none} makes a block without input (shown as not connected).

\subsection{Labels}
Labels are names, not positions: \texttt{1, 2, 5, 9} is fine, and so is a word (\texttt{stem conv 32 3x3}, then \texttt{<- stem}). The
label may also be left out: \texttt{conv 32 3x3} on its own gets the next free number. Two lines may not have the same label. The labels
are what the canvas shows in each block's header and what the layer table uses, so canvas, text and table name blocks the same way.

\subsection{Names and arithmetic: \texttt{let}}
\texttt{let} gives a value a name; every number may be an expression of names and numbers with \texttt{+ - * /}, \texttt{//}
(whole-number division), \texttt{\%} (remainder), \texttt{\^{}} (power) and parentheses:
\begin{verbatim}
let w = 32
let depth = 6
2 conv w 3x3
3 conv w * 2 3x3 stride 2
4 dense 2 ^ 10 gelu        # 1024
\end{verbatim}
Changing \texttt{w} changes every block that uses it: the quickest way to try a wider or narrower network.

\subsection{Repeating a group: \texttt{repeat}}
Identical parts are written once. A \texttt{repeat} line has a label, a count and the group in braces; the passes are chained (each
pass takes the previous pass's output), and the line's output is the last pass's last block:
\VerbatimInput[fontsize=\small,frame=single,numbers=left,numbersep=6pt,label=\file{vit\_small.nnb}]{examples/vit_small.nnb}
Inside the braces:
\begin{itemize}[nosep]
  \item blocks have \textbf{local labels} (\texttt{a}, \texttt{b}, \dots); on the canvas they become \texttt{5.1.a} \dots\
        \texttt{5.6.g} (line, pass, block);
  \item \texttt{in} is the group's input (for the first pass, the block above the \texttt{repeat} line);
  \item \texttt{i} is the pass number, 1, 2, \dots: \texttt{conv 32 * i 3x3} gets wider in every pass;
  \item the count may be a name (\texttt{repeat depth}); \texttt{repeat 12 as enc} names the group for readers.
\end{itemize}
This ViT has 49 blocks and 1\,806\,538 parameters at $3\times32\times32$; ViT-g/14 is the same pattern with \texttt{repeat 40} and
16 lines in all.

\subsection{A loop over values: \texttt{for}}
\texttt{for} repeats a group once per value, chained like \texttt{repeat}. A value can be a tuple, for several names at once:
\begin{verbatim}
5 for c in 64, 128, 256 { a conv c 3x3 }
6 for c, s in (64, 1), (128, 2), (256, 2) {
    x conv c 3x3 stride s
    y bn
  }
\end{verbatim}
\texttt{i} is also the pass number in a \texttt{for}.

\subsection{Functions: \texttt{def}}
A function is a group of blocks with values, written once and used many times. It is defined with \texttt{def}, its parameters may have
defaults, and a line calls it like a block:
\VerbatimInput[fontsize=\small,frame=single,numbers=left,numbersep=6pt,label=\file{resnet\_small.nnb}]{examples/resnet_small.nnb}
\begin{itemize}[nosep]
  \item \texttt{basic(c, s)} passes the values by position; \texttt{basic(c)} leaves \texttt{s} at its default; \texttt{basic(c=64, s=2)}
        names them.
  \item Inside the function, \texttt{in} is the input of the call (here the line above the call); the last block is the call's output.
  \item A function may call other functions (up to 16 deep: a function that calls itself is reported).
  \item Functions may be defined anywhere in the text, also below their first use.
\end{itemize}
This network has six residual blocks (66 blocks on the canvas, 791\,338 parameters): about forty lines of text.

\subsection{Several inputs, and reaching into a group}
A call (or a \texttt{repeat} / \texttt{for}) can take several inputs: \texttt{<- 3, 2.3.e} gives the call two, called \texttt{in} and
\texttt{in2} inside. A block made inside a group or a call is named by its path: \texttt{2.3.e} is line 2, pass 3, block \texttt{e} (for a
call: line, then block, as in \texttt{7.u}). Together they make U-Net's skip connections:
\VerbatimInput[fontsize=\small,frame=single,numbers=left,numbersep=6pt,label=\file{unet\_small.nnb}]{examples/unet_small.nnb}
Line 2 builds the encoder as three passes whose width depends on the pass number ($16\cdot2^{i-1}$: 16, 32, 64); lines 4--6 each take
the block above and the matching encoder output (25 blocks, 482\,050 parameters for $64\times64$ images and two classes).

\subsection{Every block}
\begin{longtable}{@{}p{3.6cm}p{5.6cm}p{5.8cm}@{}}
\toprule
\textbf{Block (other words)} & \textbf{Example} & \textbf{Settings (default)}\\
\midrule\endhead
\texttt{input} & \texttt{1 input} & none: the shape is the dataset's\\
\texttt{conv} (\texttt{conv2d}) & \texttt{conv 64 3x3 stride 2 pad 1 gelu} & filters (32), kernel \texttt{KxK} or \texttt{K} (3), \texttt{stride} (1), \texttt{pad} n / \texttt{same} / \texttt{valid} (same), \texttt{dil} (1), activation (relu), \texttt{nobias}\\
\texttt{dwconv} (\texttt{dwconv2d}) & \texttt{dwconv x2 3x3 stride 2} & multiplier \texttt{x}n (1), kernel, stride, padding, activation, \texttt{nobias}\\
\texttt{convT} (\texttt{convtranspose}) & \texttt{convT 64 2x2 stride 2 none} & filters (32), kernel (2), stride (2), \texttt{pad} (0), \texttt{outpad} (0), activation (relu)\\
\texttt{maxpool}, \texttt{avgpool} (\texttt{pool}) & \texttt{maxpool 3x3 stride 2 pad 1} & size (2), \texttt{stride} n or \texttt{size} (= size), \texttt{pad} (0)\\
\texttt{global\_pool} (\texttt{gap}) & \texttt{global\_pool max} & \texttt{avg} / \texttt{max} (avg)\\
\texttt{upsample} & \texttt{upsample x2 bilinear} & scale \texttt{x}n (2), \texttt{nearest} / \texttt{bilinear}\\
\texttt{bn}, \texttt{ln}, \texttt{gn} & \texttt{gn 8} & GroupNorm: groups (8; must divide the channels)\\
\texttt{dropout} & \texttt{dropout 0.3} & rate (0.25)\\
\texttt{flatten} & \texttt{flatten} & none\\
\texttt{dense} (\texttt{linear}) & \texttt{dense 128 gelu} & units (128), activation (relu)\\
\texttt{relu}, \texttt{gelu}, \texttt{silu}, \texttt{tanh}, \texttt{sigmoid}, \texttt{leaky\_relu} & \texttt{gelu} & an activation block\\
\texttt{add}, \texttt{concat} (\texttt{cat}), \texttt{mul} (\texttt{multiply}) & \texttt{add <- 2, 5} & two or more inputs; Concat joins channels in the order given\\
\texttt{to\_tokens} & \texttt{to\_tokens patches} & \texttt{patches} (every pixel a token) / \texttt{rows}\\
\texttt{cls\_pos} & \texttt{cls\_pos} & class token + position embedding; \texttt{nocls}: positions only\\
\texttt{attention} (\texttt{attn}) & \texttt{attention heads 8 drop 0.1} & \texttt{heads} (4; must divide the token width), \texttt{drop} (0)\\
\texttt{token\_pool} & \texttt{token\_pool mean} & \texttt{cls} / \texttt{mean} (cls)\\
\texttt{output}, \texttt{seg\_output} & \texttt{output} & none: classes from the dataset\\
\bottomrule
\end{longtable}
The activation words after a \texttt{conv}, \texttt{dwconv}, \texttt{convT} or \texttt{dense} set its activation (\texttt{none}: without);
on their own line they are an Activation block. Block words may be written in capitals (\texttt{Conv}); the text the application prints
uses the short lower-case words.

\subsection{The rules in brief}
\begin{verbatim}
net "name"                                   name of the design (optional)
let NAME = EXPR                              a named value
def NAME(P, Q = EXPR) { LINES }              a function (P, Q: parameters; Q has a default)
LABEL BLOCK SETTINGS [<- REF, REF]           one block; without <- it takes the line above
LABEL NAME(EXPR, Q=EXPR) [<- REF, REF]       a function call (inside: in, in2, ...)
LABEL repeat EXPR [as NAME] [<- REF] { LINES }       passes chained; i = pass number
LABEL for A, B in (EXPR, EXPR), ... [<- REF] { LINES }
REF:  a label | in | in2 | a path like 5.2.b    EXPR: numbers, names, + - * / // % ^ ( )
\end{verbatim}

\subsection{Mistakes and what the application says}
A text with a problem is not applied: the canvas keeps the last valid design, the status says ``not applied'', the line number in the
gutter turns red and the message below the text names the line (click it to jump there). Typical messages:
\begin{longtable}{@{}p{4.6cm}p{10.4cm}@{}}
\toprule
\textbf{Written} & \textbf{Message}\\
\midrule\endhead
\texttt{conv 64 3x3 stide 2} & 'stide' means nothing for conv (did you mean 'stride'?)\\
\texttt{convv 64 3x3} & unknown block 'convv' (did you mean 'conv'?)\\
\texttt{add <- 2, 7} (no line 7 above) & '7': there is no block '7' above this line (a block may only use blocks above it)\\
two lines labelled \texttt{2} & the label '2' is used twice\\
\texttt{dense 10 pad 1} & 'pad' means nothing for dense\\
\texttt{conv 8 3x3 filtrs=4} & conv has no setting 'filtrs' (its settings: activation, bias, dilation, filters, kernel, padding, stride)\\
\texttt{conv 8 3x3 99} & what is the number 99? write it as name=value (e.g.\ filters=99)\\
a \texttt{\{} without \texttt{\}} & a '\{' is not closed with '\}'\\
\texttt{f(1, 2)} for \texttt{def f(c)} & 'f' takes 1 value(s), more were given\\
\texttt{repeat k} without \texttt{let k} & repeat: 'k' is not a number\\
\texttt{repeat 5000 \{ \dots\ \}} too large & 5000 passes × 7 blocks = 35000 blocks, more than the 13000 a design may have (Settings → Design size \dots). At most 1856 passes fit here.\\
\bottomrule
\end{longtable}
Shape problems (an Add of different shapes, a kernel larger than the image) are found by the same checks as on the canvas and listed under
the text with the line of their block; click \emph{block 5.2.a} to see it selected on the canvas. \kbd{Cmd}-click on any label in the text
also selects its block.

\subsection{How the text and the canvas stay in step}
\begin{itemize}[nosep]
  \item \textbf{Typing:} a quarter of a second after the last key, a valid text replaces the design (one undo step). Blocks keep their
        place on the canvas when their label and kind stay the same; new blocks are placed next to their input.
  \item \textbf{Moving blocks} on the canvas keeps the text exactly as written.
  \item \textbf{Changing a setting} of a block that has its own line rewrites only that line (its label, inputs and comment stay).
  \item \textbf{Adding, removing or connecting} blocks of plain lines changes only those lines: a new block gets the next free number and
        goes after its input; a line that took ``the line above'' gets an explicit \texttt{<-} when that changed; a block inserted in front
        of a function call or loop becomes its input. Names, functions, loops and comments stay.
  \item \textbf{Changing a block that a function or loop made} would change every call, so the text is printed again from the blocks:
        identical consecutive groups become \texttt{repeat}, the rest plain lines. \kbd{Cmd+Z} brings the hand-written text back.
  \item \ui{Format} prints the text again from the blocks on purpose; \ui{Copy} copies it.
\end{itemize}
In the layer table, the blocks of one \texttt{repeat}, \texttt{for} or call line form one group named after the line.

\subsection{A design that uses everything}
\file{doc/examples/everything.nnb} in the repository is a hybrid CNN and Transformer written only as text: functions calling functions
with defaults and named values, a function with two inputs, \texttt{for} with tuples, \texttt{repeat} with \texttt{i}, \texttt{let} and
every operator, paths into loop passes, a left-out label, comments, and every block kind except the segmentation output (89 blocks,
186\,086 parameters, no warning at $3\times32\times32$). The test suite builds it and every example of this chapter, and GitHub Actions
compares its layers with real PyTorch.
